\scanpage{35}
Posant
\[
 \unsure{\xi|\zeta}^{*}_{E_{\unsure{y|\eta}}}(t)=t^{\chi^*(E_y)/2}L^*_{E_y}(t)
\]
on trouve
\[
 \boxed{\xi^*_{E_y}\left(\frac{1}{qt}\right)=\varepsilon'{}^*(E_y)\xi^*_{E_y}(t)}
\]
\[
 \left\{\begin{aligned}
 \varepsilon'{}^*(E_y)&=\delta^*(E_y)q^{-\chi^*(E_y)/2}(-1)^{\chi^*(E_y)}\\
 q&=p^{1+\rho}
 \end{aligned}\right.
 \qquad\text{\unsure{est|ou|où} \unsure{une|avec} \unsure{cte|etc}}
\]

\uline{Un reproche à faire à} $L^*_{E_y}(t)$ : \uline{ce n'est pas additif} en $E$. De façon précise, si on a
\diagram{\begin{tikzcd}[ampersand replacement=\&]
0\arrow[r] \& E'_y\arrow[r,no head] \& E_y\arrow[r,no head] \& E''_y\arrow[r] \& 0
\end{tikzcd}}
On en \deleted{\unsure{en|aura}} conclut
\diagram{\begin{tikzcd}[ampersand replacement=\&]
0\arrow[r,no head] \& i_*(E'_y)\arrow[r,no head] \& i_*(E_y)\arrow[r] \& i_y^{\unsure{*|2}}(E''_y)\arrow[r] \& Q\arrow[r] \& 0
\end{tikzcd}}
où $Q\deleted{=\operatorname{Ker}\left(R^1i_*(E'_y)\mathbin{-}R^1i_*(E_y)\right)}$
\inserted{$\simeq\coprod_{\substack{x\ \text{pt ramifié}\\\text{de }E''_y}}j_x\unsure{(\quad)|!}$}
est \unsure{concentré} \deleted{\unsure{aux}} \unsure{ès|les} points du \unsure{lieu} de \textquotedblleft ramification\textquotedblright\ de $E''_y$.

D'où
\[
 \left\{\begin{aligned}
 L^*_{E_y}(t)&=L^*_{E'_y}(t)L^*_{E''_y}(t)L_Q(t)^{-1}
 =L^*_{E'_y}(t)L^*_{E''_y}(t)\prod_x P_{Q_x}(t)\\
 Q_x&=\operatorname{Coker}\left(E_x\mathbin{-}E''_x\right)
 \end{aligned}\right.
\]
\margin{$H^2(\bar X,\bar E)={}$ dual de $H^0(\bar X,\bar E\deleted{\unsure{(1)|(\rho)}})(\unsure{\rho+1|\rho-1})$}
% archive sheet number: 34
